% ============================================================
% Chapter 5 -- Results and Analysis
% ============================================================
% Writing-suggestions entries applied: C5-01 (calib/purity split),
%   C5-MMLU-01 (dual forgetting instruments), C5-12 (existing McNemar +
%   bootstrap helpers), GEN-12 (significance required), GEN-03
%   (projected post-calibration weights).
%
% Phase 4m.2: This is the first committed content for Chapter 5 --
% the evaluation-protocol preface for Section~5.1 that defines the
% seven metric tables (5.1 Headline, 5.2 Calibration, 5.3 Hallucination
% Decomposition, 5.4 Grounding, 5.5 Data Purity, 5.6 Continual
% Learning, 5.7 Statistical Significance) before any result is
% reported. Subsequent sections (§§5.2--5.9) are populated from live
% experiment outputs and are not drafted in this phase.
%
% Note: citation keys lopezpaz2017gradient (Lopez-Paz \& Ranzato 2017,
% NeurIPS, ``Gradient Episodic Memory for Continual Learning''),
% brier1950verification (Brier 1950, Monthly Weather Review), and
% mcnemar1947note (McNemar 1947, Psychometrika) are now present in
% bibliography/references.bib alongside the 2026-04-17 Ch~2 benchmark
% additions (joshi-etal-2017-triviaqa, kwiatkowski-etal-2019-natural,
% li-etal-2023-halueval, clark2018think, hendrycks2021measuring).
% ============================================================

% ---------------------------------------------------------------
\section{Experimental Setup}
\label{sec:exp-setup}
% ---------------------------------------------------------------

\Cref{sec:exp-setup} defines the evaluation protocol used across the
remainder of this chapter: the seven benchmarks \caem{} is evaluated on,
the metric suite applied to each benchmark at each cycle, and the eight
external baselines against which \caem{} is compared (B1--B8), plus an
optional STaR ceiling reference reported when compute permits. The
metric suite is deliberately broader than prior single-number
evaluations of hallucination systems: a single exact-match score cannot
disentangle accuracy from calibration, grounding, confabulation, or
forgetting, and a thesis claiming an \emph{architectural} intervention
must demonstrate behaviour on each of these axes separately before any
aggregate claim is accepted.

% ---------------------------------------------------------------
\subsection{Benchmarks}
\label{sec:benchmarks}
% ---------------------------------------------------------------

\caem{} is evaluated on seven factual-QA benchmarks spanning three
in-domain training benchmarks (FEVER~\cite{thorne-etal-2018-fever},
TriviaQA~\cite{joshi-etal-2017-triviaqa}, and Natural
Questions~\cite{kwiatkowski-etal-2019-natural}) and four out-of-domain
transfer benchmarks
(HaluEval~\cite{li-etal-2023-halueval},
TruthfulQA~\cite{lin-etal-2022-truthfulqa},
StrategyQA~\cite{10.1162/tacl_a_00370},
and ARC-Challenge~\cite{clark2018think}),
with MMLU~\cite{hendrycks2021measuring} held out as the out-of-distribution
retention control for the forgetting-tolerance guard.
HaluEval is included specifically to test the confident-error axis
of \caem{}'s claim: the benchmark carries per-sample hallucination
labels, so the stored-set correctness rate can be cross-checked
against an external hallucination judgement rather than relying
purely on EM.
This split cleanly separates the in-domain contribution of each
architectural mechanism from its zero-shot transfer effect; the
ablation study in~\cref{sec:ablation} exploits this separation.
Benchmark selection rationale and dataset statistics are given in
\cref{tab:benchmark-overview}.

\begin{table}[t]
\centering
\caption{Benchmark panel overview. Split source, sample count per
cycle, answer format, and metric used for each of the seven
factual-QA benchmarks plus the MMLU retention control.}
\label{tab:benchmark-overview}
\begin{tabular}{lllll}
\toprule
Benchmark & Split & $N$ & Answer format & Primary metric \\
\midrule
FEVER               & paper\_dev & 500 & 3-way label    & Label accuracy \\
TriviaQA            & validation & 500 & Free-form      & EM / F1 \\
Natural Questions   & validation & 500 & Span           & EM / F1 \\
HaluEval            & QA split   & 500 & Binary         & Hallucination F1 \\
TruthfulQA          & generation & 500 & Free-form      & ROUGE-L $\ge 0.15$ \\
StrategyQA          & dev        & 500 & Binary         & Label accuracy \\
ARC-Challenge       & test       & 500 & 4-way MCQ      & Label accuracy \\
MMLU (retention)    & 57 subjects& 200 & 4-way MCQ      & Retention ratio $\rho$ \\
\bottomrule
\end{tabular}
\end{table}

% ---------------------------------------------------------------
\subsection{Metrics}
\label{sec:metrics}
% ---------------------------------------------------------------

Every cycle of every benchmark produces a per-sample record containing
the generated answer, the reasoning chain, the routing tier, wall-clock
latency, and the full set of verification signals emitted by the
\textsc{UnifiedVerifier} at Stage~5. The metrics below are computed
from these per-sample records and are reported in seven dedicated
tables (Tables~\ref{tab:headline}--\ref{tab:sig-test}). We organise the
suite into seven categories, each answering a distinct evaluation
question.

\paragraph{Accuracy and latency (Table~\ref{tab:headline}).}
Primary correctness is reported with exact match (EM) and token-level
F1 on free-form answers (TriviaQA, NQ), label accuracy on classification
benchmarks (FEVER three-way, ARC-Challenge four-way,
StrategyQA binary), and ROUGE-L$\geq 0.15$ on TruthfulQA following
Lin et al.~\cite{lin-etal-2022-truthfulqa}.
Because \caem{}'s architectural value lies in its routing selectivity,
each benchmark's EM is also reported
\emph{stratified by routing tier} -- $\mathrm{EM}_{T_1},
\mathrm{EM}_{T_2}, \mathrm{EM}_{T_3}$ -- so that a rise in Tier~1
fraction (memory hits) across cycles can be read alongside per-tier
correctness rather than being averaged away.
Wall-clock latency is reported per tier for the same reason: Tier~1
skips full decoding, so the efficiency claim reduces to ``Tier~1
latency is substantially below Tier~2 and Tier~3 latency, and Tier~1
share grows across cycles.'' Both quantities are computed via
\texttt{em\_by\_tier} and \texttt{mean\_latency\_by\_tier} in
\texttt{eval/metrics.py}.

\paragraph{Calibration (Table~\ref{tab:calibration}).}
Calibration is quantified with three complementary metrics computed on
the pooled stream of $(\ustored, \mathrm{correct})$ pairs from every
stored episode across all seven benchmarks at each cycle.
The \emph{Expected Calibration Error} (ECE; Guo et
al.~\cite{pmlr-v70-guo17a}) measures average absolute gap between
confidence and accuracy over $M=10$ equal-width bins.
The Brier score~\cite{brier1950verification}
complements ECE: it is a proper scoring rule that penalises both
miscalibration and low resolution, and it does not require bin
aggregation -- a system that is badly calibrated in ways that
counter-balance inside bins will still receive a poor Brier score
even as ECE remains low.
The \emph{area under the ROC curve} (AUROC) treats $\ustored$ as a
predictor of correctness and measures rank discrimination independent
of scale; a high AUROC with poor ECE indicates that the confidence
signal is useful for routing even before calibration.
Bin-wise \emph{reliability diagrams} (Figure~\ref{fig:reliability})
ground these scalars visually using the 10-bin decomposition returned
by \texttt{reliability\_bins}.

\paragraph{Hallucination decomposition (Table~\ref{tab:halluc}).}
A single hallucination-rate scalar conflates distinct failure modes,
so we report three rates side by side.
The \emph{stored-hallucination rate} is the fraction of stored
episodes that have $\mathrm{EM}=0$ and
$\ustored \geq 0.50$ -- i.e., stored-with-confidence errors, which is
the quantity the data purity theorem bounds.
The \emph{internal-confidence confabulation rate} is the fraction
with $\mathrm{EM}=0$ and $u_\text{internal} \geq 0.70$,
capturing over-confident generation regardless of whether the episode
was stored; this follows the hidden-state-based notion of confabulation
used by Farquhar et al.~\cite{farquhar2024semantic}.
The \emph{semantic-entropy confabulation rate} applies the same
threshold on $(1 - h_{\text{norm}})$, where $h_{\text{norm}}$ is the
normalised semantic entropy of \caem{}'s $K=10$ answer cluster, serving
as a direct Farquhar replication on \caem{}'s own samples.
Reporting all three prevents the case where reducing one form of
hallucination silently inflates another. All variants are computed
by the generic helper \texttt{confabulation\_rate} with different
(signal, threshold) pairs.

\paragraph{Grounding (Table~\ref{tab:grounding}).}
Two grounding diagnostics isolate the retrieval pathway.
\emph{Unsupported-correct rate} is the fraction of
$\mathrm{EM}=1$ samples with $p_{\text{ground,max}} < 0.30$ -- correct
answers that retrieval evidence did not support.
\emph{Ungrounded-assertion rate} is the fraction of all samples with
$p_{\text{ground,max}} \leq 0.30$ and $\ustored \geq 0.50$ -- the
evaluation-side mirror of the Stage~5 early-exit confabulation gate
($u_\text{internal} \geq 0.70$ and
$p_{\text{ground,max}} \leq 0.20$) generalised so the same definition
serves both the online gate in~\cref{sec:system-overview} and
this offline diagnostic. The first quantifies parametric-memory
reliance; the second quantifies confidently-unsupported output
regardless of correctness.

\paragraph{Data purity (Table~\ref{tab:purity}).}
The verifier-decision breakdown -- counts and fractions of
\textsc{STORE}, \textsc{DEFER}, \textsc{ABSTAIN}, and
\textsc{DISCARD} decisions per cycle -- is reported alongside the
empirical verifier true-positive rate $\alpha$, the base generation
accuracy $p$, the theorem's predicted post-storage purity
$P_\text{theory} = p\alpha / (p\alpha + (1-p)(1-\alpha))$, and the
measured purity $P_\text{obs}$. This table is the empirical
instantiation of the purity theorem from~\cref{sec:purity-theorem};
all five quantities are computed from the same purity-validation split
defined in~\cref{sec:benchmarks}. The \textsc{STORE} denominator is
obtained from \texttt{decision\_breakdown}.

\paragraph{Continual learning (Table~\ref{tab:continual}).}
Forgetting is reported with two instruments, following the dual-metric
design in~\cref{sec:self-improvement}. \emph{MMLU retention} is the
ratio of Cycle-$n$ MMLU accuracy (200 questions, shuffled seed~42
across the 57 subjects) to the Cycle-$0$ pristine-model baseline; it
is domain-neutral and never trained on, and is the headline forgetting
metric. In addition, we report \emph{backward transfer} and
\emph{forward transfer} following the gradient-episodic-memory
convention of Lopez-Paz and
Ranzato~\cite{lopezpaz2017gradient}.
Backward transfer averages $R_{T,i} - R_{i,i}$ over earlier
benchmarks $i$; positive values indicate that later cycles improved
earlier benchmarks, and negative values indicate catastrophic
forgetting. Forward transfer averages $R_{i-1,i} - \bar b_i$ relative
to the Cycle-$0$ baseline and measures whether learning on earlier
cycles helps benchmarks not yet reached.
Reporting both in addition to MMLU retention answers
the per-benchmark-per-cycle drift question that MMLU retention (being
a single scalar) cannot address.

\paragraph{Unified efficacy (CAEM Efficacy Score).}
Because a reviewer must read the seven tables above quickly, we report
a single composite scalar -- the \emph{\caem{} Efficacy Score} (CES) --
as the geometric mean of five orthogonal axes:
\[
\mathrm{CES} = \bigl(\mathrm{ACC} \cdot \mathrm{EPI} \cdot \mathrm{RET}
                    \cdot \mathrm{CAL} \cdot \mathrm{VER}\bigr)^{1/5},
\]
with $\mathrm{ACC}$ the mean EM across all seven benchmarks,
$\mathrm{EPI} = 1 - \text{(pooled confident-error rate)}$,
$\mathrm{RET} = \min(\text{MMLU retention ratio},\,1)$,
$\mathrm{CAL} = 1 - 2\min(\mathrm{ECE},\,0.5)$,
and $\mathrm{VER} = \max(2\beta - 1,\,0)$ where $\beta$ is the
verifier's balanced accuracy on labelled validation data. Each
component is clamped to $[0.01, 1]$ so that a transient failure on
one axis does not zero the score, but the geometric form ensures
that a genuine collapse on any axis (for instance, catastrophic
forgetting destroying $\mathrm{RET}$) pulls CES down rather than
being averaged away. CES is a thesis-introduced aggregate used
\emph{only} as a visualisation and headline scalar; the seven
component tables remain the authoritative quantitative record, and
CES is always reported with all five constituent values.

\paragraph{Statistical significance (Table~\ref{tab:sig-test}).}
All accuracy comparisons between \caem{} and baselines are tested
for significance using McNemar's
test~\cite{mcnemar1947note} with the Edwards continuity correction,
applied to paired per-question correct/incorrect outcomes. Non-parametric
$95\%$ confidence intervals for EM are computed by bootstrap resampling
with $B = 1000$ resamples. Both computations use the existing
\texttt{mcnemar\_test} and \texttt{bootstrap\_ci} helpers in
\texttt{eval/metrics.py}, which were written and tested before the
experiment phase.

Because \caem{} is compared against eight baselines across seven
benchmarks, the raw test count is $8 \times 7 = 56$ pairwise McNemar
tests per cycle, and a naive $p < 0.05$ threshold would inflate the
family-wise error rate to approximately $1 - 0.95^{56} \approx 0.94$.
We therefore apply the Holm--Bonferroni step-down
correction~\cite{holm1979simple} within each benchmark family: the
seven per-benchmark \caem{}-vs-baseline $p$-values are sorted in
ascending order, and the $k$-th smallest is compared against
$0.05 / (m - k + 1)$ where $m = 8$ is the number of baselines. This
is uniformly more powerful than plain Bonferroni while preserving the
family-wise error rate at $\alpha_{\text{fam}} = 0.05$, and is the
accepted correction for structured pairwise-comparison panels of
this size. A result is marked ``significant'' in
Table~\ref{tab:sig-test} only when it survives this step-down
correction; results that are directionally positive but fail the
correction are reported with both the raw and corrected $p$-values
so the reader can distinguish a suggestive trend from a confirmed
effect.

We separately require a minimum practical-significance effect size of
$\Delta \mathrm{EM} \geq 0.02$ (absolute, two percentage points on a
$500$-sample benchmark). McNemar's statistical test has high power at
this sample size and will flag effects as small as $\Delta \approx 0.01$
as significant, but differences below two points are within the cycle-
to-cycle variation \caem{} itself exhibits on a held-out seed, so we
do not claim architectural advantage below that threshold even when
the test is significant. Table~\ref{tab:sig-test} therefore reports
two markers: $\dagger$ for Holm-corrected $p < 0.05$, and $\star$ for
the conjunction of Holm-corrected $p < 0.05$ \emph{and}
$\Delta \mathrm{EM} \geq 0.02$. The thesis's architectural claims are
licensed only by the $\star$ marker; the $\dagger$ marker identifies
directionally positive effects whose magnitude is too small for the
claim to be practically meaningful, even though the underlying
comparison is statistically sound.

Where \caem{} loses a comparison (a baseline's EM exceeds \caem{}'s on
some benchmark), we apply the same two-stage rule in reverse: a
baseline is reported as strictly superior to \caem{} on a benchmark
only if $\Delta \mathrm{EM} \leq -0.02$ with Holm-corrected
$p < 0.05$. This symmetric framing prevents the write-up from silently
suppressing the baseline-wins case, which is a common failure mode
in architectural evaluations and one that the ablation
discussion~(\S\ref{sec:ablation-results}) treats as legitimate
evidence rather than as noise to be explained away.

\paragraph{Summary.}
The seven metric categories above answer seven distinct evaluation
questions: \emph{Is \caem{} more accurate and faster than its
baselines?} (Table~\ref{tab:headline});
\emph{Are \caem{}'s confidence signals well-calibrated?}
(Table~\ref{tab:calibration});
\emph{Has \caem{} reduced confident errors across every hallucination
sub-type?} (Table~\ref{tab:halluc});
\emph{Is \caem{}'s retrieval pathway earning its slot?}
(Table~\ref{tab:grounding});
\emph{Does the purity theorem's prediction match measurement?}
(Table~\ref{tab:purity});
\emph{Does self-improvement proceed without catastrophic forgetting
in any direction?} (Table~\ref{tab:continual});
\emph{Are the claimed improvements statistically significant?}
(Table~\ref{tab:sig-test}).
CES provides the single-scalar visualisation of these axes.
The remainder of this chapter populates each table in turn.

% ---------------------------------------------------------------
\subsection{Baselines}
\label{sec:baselines}
% ---------------------------------------------------------------

\caem{} is compared against eight external baselines spanning the
three families of prior systems identified in
\cref{sec:external-baselines}: prompt-level reasoning, retrieval
augmentation, and training-time self-improvement. All eight share the
\mbox{Flan-T5-Large} (780M) backbone wherever a backbone swap would not
invalidate the method, and all eight are run on the same seven-benchmark
factual-QA panel with the same 500-sample draw per benchmark, so any
difference in score isolates the contribution of the architectural
component that distinguishes the baseline from \caem{}.

\begin{description}
  \item[B1 --- Zero-shot Flan-T5-Large.]
  The same 780M-parameter backbone \caem{} fine-tunes, invoked with no
  CoT prefix, no retrieval, and no verification. This is the empirical
  floor: any gain \caem{} reports must exceed this score, otherwise the
  gain is attributable to model capacity alone.

  \item[B2 --- Chain-of-Thought.]
  Flan-T5-Large with the ``Let's think step by step'' prefix
  \cite{wei2022chain}, no retrieval, no verification. Isolates the
  contribution of intermediate reasoning on the same backbone.

  \item[B3 --- DPR-RAG.]
  Dense Passage Retrieval \cite{karpukhin-etal-2020-dense} with
  $k=5$ Wikipedia passages concatenated into a 384-token context and
  generated from the same Flan-T5-Large backbone
  \cite{NEURIPS2020_6b493230}. Isolates the contribution of external
  grounding without any verification or memory.

  \item[B4 --- CoT + DPR-RAG.]
  B3 with the CoT prefix prepended. Isolates the contribution of
  reasoning over retrieved evidence.

  \item[B5 --- FLARE.]
  Active retrieval-augmented generation \cite{jiang-etal-2023-active}
  with confidence floor $\theta = 0.4$ and a 64-token look-ahead
  window, using the same DPR index as B3 and B4. Isolates the
  contribution of confidence-gated retrieval, which is the closest
  prior mechanism to \caem{}'s tier routing; the B5-versus-\caem{}
  gap measures whether three tiers and the nine-signal composite
  meaningfully improve on a single-tier confidence gate.

  \item[B6 --- Vanilla fine-tuning.]
  Flan-T5-Large trained for the same ten cycles as \caem{} on the
  verified episode pool, but without the $L_2$ anchor, without the
  MMLU retention guard, and without the episodic memory or the
  verifier at inference time, so every selected episode is treated as
  a clean label. Isolates the contribution of \caem{}'s
  retention-protection machinery.

  \item[B7 --- EWC-only fine-tuning.]
  B6 with the $L_2$ anchor and the MMLU retention guard added
  \cite{doi:10.1073/pnas.1611835114} but with the verifier, the
  nine-signal composite, the memory, and the routing all still
  stripped. Isolates the contribution of the
  verification-and-memory machinery once retention machinery is held
  equal.

  \item[B8 --- Self-RAG (citation-only).]
  Asai et al.~\cite{asai2024selfrag} is the most ambitious
  retrieve-and-critique pipeline in the prior literature but its
  critic tokens are tied to a LLaMA-family backbone. Re-targeting
  them to Flan-T5-Large would change enough of the system that the
  resulting number would measure backbone and prompt-engineering
  differences more than architectural differences, so Self-RAG is
  included as a positioning reference rather than as a replicated
  numerical baseline.
\end{description}

Optional ceiling reference: \textbf{STaR}
\cite{NEURIPS2022_639a9a17}, the pure iterative-refinement scheme
reviewed in \cref{sec:external-baselines}, is reported separately where
compute permits, and serves as an aspirational upper bound for what a
verification-free bootstrap can reach on the same backbone.

Each baseline lacks exactly one of \caem{}'s architectural components;
the pairwise comparisons in \cref{sec:main-results} isolate the
contribution of that component. Implementation details, shared backbone
model, and compute budget are given in~\cref{sec:impl-details}.

% ---------------------------------------------------------------
\subsection{Implementation Details}
\label{sec:impl-details}
% ---------------------------------------------------------------

All results in this chapter are produced by a single implementation
run against a pinned set of dependencies and a fixed hardware
envelope. This subsection records the substrate so that any number
reported in \S\ref{sec:main-results}--\S\ref{sec:ablation-results} is
traceable to a specific model checkpoint, GPU, configuration file,
and random seed.

\paragraph{Backbone and supporting models.}
The generator is Flan-T5-Large~(\cite{JMLR:v25:23-0870}, $780$M
parameters), loaded at \texttt{bfloat16} precision. Sentence-level
embeddings for episodic memory keys, retrieval similarity, and
semantic-entropy clustering are produced by the frozen
\texttt{all-mpnet-base-v2} sentence
encoder~(\cite{reimers-gurevych-2019-sentence}, $110$M parameters,
$768$-dim output). NLI entailment and contradiction probabilities
feeding the verifier's nine-signal composite are produced by the
frozen \texttt{roberta-large-mnli} cross-encoder~(\cite{liu2019roberta},
$355$M parameters). The same three checkpoints are used by every
\caem{} variant and by every retrieval-family baseline (B3 DPR-RAG,
B4 CoT+RAG, B5 FLARE); backbone identity is held constant across the
comparison panel so that score differences isolate architectural
contributions rather than capacity differences.

\paragraph{Hardware.}
Cycle-level training and evaluation run on a single NVIDIA
RTX~4090 ($24$~GB VRAM) rented through Vast.ai, paired with a
Ryzen-class host CPU and $64$~GB of system RAM. The L$_2$ anchor
$\theta_{\text{prev}}$ co-resides on GPU with the training parameters
so that the anchor penalty can be computed without host-device
transfers each step, which doubles the effective weight footprint
during a cycle and motivates the \texttt{bfloat16} precision choice.
The episodic-memory FAISS index stays on GPU for vector search and
spills to CPU only for periodic save; the Wikipedia DPR passage index
($21$M vectors) is held on CPU throughout because its size exceeds
what can be co-resident with the generator on a single $24$~GB card.

\paragraph{Cycle structure and sample budgets.}
\caem{} is trained for $N = 10$ self-improvement cycles following the
schedule in \S\ref{sec:self-improvement}. At every cycle the
generator is evaluated on the full seven-benchmark panel with
$n_{\text{eval}} = 500$ questions per benchmark drawn under the
fixed split assignments documented in Table~\ref{tab:benchmark-overview}.
Four additional disjoint sample splits support the per-cycle
diagnostics: $n_{\text{cal}} = 500$ for temperature calibration and
reliability-diagram fitting, $n_{\text{purity}} = 500$ for the
verifier-balanced-accuracy estimate and the
$\lvert P_{\text{theory}} - P_{\text{obs}} \rvert$ comparison that
populates Table~\ref{tab:purity}, $n_{\text{MMLU}} = 200$ shuffled
across the 57 MMLU subjects for the retention guard, and
$n_{\text{screen}} = 1\,500$ for the Phase~1 screening sweep that
ranks the sixteen training-time ablation variants before the
confirmatory panel. All five splits are non-overlapping within a
benchmark and the split assignment is fixed across cycles and across
all ablation variants.

\paragraph{Cold-start memory.}
Cycle~$0$ begins with a pre-populated episodic store of $447$
verified episodes drawn from an 80/20 stratification of the three
in-domain training benchmarks (FEVER, TriviaQA, Natural Questions)
under the same Stage~5 verifier used online, seeded with the pristine
Flan-T5-Large backbone. This is documented in detail in the
cold-start section of \S\ref{sec:self-improvement}; the role of the
$447$ episodes is to give Tier~1 retrieval a non-empty index on the
very first query of Cycle~$0$ so that the routing decision can be
exercised from the first sample rather than starting with an
inevitable Tier~2/Tier~3 fallback.

\paragraph{Reproducibility.}
Every cycle run is controlled by a single configuration dataclass
(\texttt{caem.config.CAEMConfig}) that records every numeric constant
in the system; each constant carries a literature~[LIT], design~[DES],
or calibration~[CAL] tag per~\S\ref{sec:implementation}, and the
dataclass asserts at load that no field is missing. Random seeds are
pinned at three levels: Python's \texttt{random} (\texttt{seed=42}),
NumPy (\texttt{np.random.seed(42)}), and PyTorch
(\texttt{torch.manual\_seed(42)}, \texttt{torch.cuda.manual\_seed\_all(42)}),
with \texttt{torch.backends.cudnn.deterministic = True} for
convolutional kernels. Benchmark splits are materialised once to
disk with a frozen sample ordering so that re-running Cycle-$n$ months
later under the same seed produces the same per-sample records. The
framework versions that any reported number is bound to are Python
$3.10$, PyTorch $2.1$, HuggingFace Transformers $4.36$, Datasets
$2.16$, sentence-transformers $2.3$, and \texttt{faiss-cpu} $1.7$.

\paragraph{Compute budget summary.}
One complete \caem{} run (Cycles~$0$--$10$ across the seven benchmarks
plus MMLU retention, plus calibration and purity diagnostics) occupies
approximately $18$ GPU-hours on the RTX~4090 envelope. The external
baseline panel (\S\ref{sec:baselines}, B1--B7) adds approximately
$16$ GPU-hours sequentially ($\sim 3.5$ hours for the five inference
baselines B1--B5 combined, $\sim 12$ hours for the two training
baselines B6--B7 combined, each of which runs for the same $10$
cycles as \caem{}). The optional STaR ceiling reference adds a further
$\sim 8$ GPU-hours when compute permits; Self-RAG is included as a
positioning reference rather than a replicated numerical baseline
(\S\ref{sec:baselines}, B8) and contributes no GPU time. The Phase~1
screening sweep over the sixteen training-time variants at three
cycles and $n_{\text{screen}} = 1\,500$ occupies a further $\sim 24$
GPU-hours before the confirmatory panel narrows to the subset that
earns full $10$-cycle runs.

% ---------------------------------------------------------------
\section{Ablation Methodology}
\label{sec:ablation}
% ---------------------------------------------------------------

This section describes the ablation protocol that populates the
ablation table in \S\ref{sec:ablation-results}.

\paragraph{Training-time versus inference-time variants.}
\caem{} is a stateful system whose weights evolve across cycles.
An ablation that merely disables a mechanism at inference --- while
evaluating on weights that were already \emph{shaped} by that mechanism
in previous cycles --- does not measure the mechanism's contribution;
it measures a counterfactual confound. We therefore classify every
ablation variant as either \emph{inference-time} (the mechanism can be
flipped off without changing training history, so the reference-run
weights may be reused) or \emph{training-time} (the mechanism alters
storage admission, verification decisions, the self-improvement loss,
or any other signal that feeds back into subsequent cycles, so the
variant requires its own full cyclic rerun from the pristine baseline).
Under strict application of this principle, only two variants are
inference-time: the full system and
\texttt{no\_self\_improvement}; all 14 remaining variants are
training-time and obtain their own 10-cycle trajectory.

\paragraph{Variant registry.}
A named registry of 16 variants is defined in
\texttt{caem/ablation/variants.py} and dispatched by
\texttt{scripts/run\_cyclic\_ablation.py}. Each cyclic rerun writes its
cycle-wise CES axes to
\texttt{outputs/ablation/<variant>/seed\_<N>/ces\_axes\_per\_cycle.json},
and \texttt{scripts/aggregate\_ablation.py} assembles the per-variant
rows for the ablation table. Table~\ref{tab:ablation-registry}
summarises the registry.

\begin{center}
\captionof{table}[Ablation Registry]{Sixteen named ablation variants.
Training-time variants each receive their own cyclic rerun; inference-time
variants reuse the reference-run weights.}
\label{tab:ablation-registry}
\renewcommand{\arraystretch}{1.2}
\footnotesize
\begin{tabular}{@{}p{4.2cm} p{1.8cm} p{8.2cm}@{}}
\toprule
\textbf{Variant} & \textbf{Class} & \textbf{What it disables / forces} \\
\midrule
\texttt{full} & inference & Reference run; no mechanism disabled. \\
\texttt{no\_self\_improvement} & inference & Disable Stage~8 fine-tuning; freeze weights at $\theta_0$. \\
\texttt{no\_tier1} & training & Force Tier~2/3 for every query; memory never serves answers. \\
\texttt{no\_retrieval} & training & Disable Tier~3 RAG retrieval entirely. \\
\texttt{no\_verification} & training & Accept every generation into memory regardless of $\ustored$. \\
\texttt{no\_nli} & training & Drop the NLI entailment signal from $\ustored$. \\
\texttt{no\_selfcons} & training & Drop the self-consistency signal from $\ustored$. \\
\texttt{no\_entropy} & training & Drop the semantic-entropy signal from $\ustored$. \\
\texttt{no\_routing} & training & Disable OR-condition safety veto; route purely on $u_{\text{pre}}$. \\
\texttt{no\_cold\_start} & training & Begin Cycle~0 with empty memory (skip 447-episode seed). \\
\texttt{no\_abort\_guard} & training & Disable $\rho < 0.93$ forgetting circuit-breaker. \\
\texttt{no\_l2\_anchor} & training & Set $\lambda = 0$; fine-tune without L2 anchor toward $\theta_{\text{prev}}$. \\
\texttt{no\_retroverify} & training & Skip end-of-cycle retroactive re-verification pass. \\
\texttt{no\_memory\_prune} & training & Disable memory pruning; allow unbounded store growth. \\
\texttt{no\_mcdropout} & training & Drop MC-Dropout signal from $u_{\text{pre}}$. \\
\texttt{lower\_u\_threshold} & training & Replace $\ustored \geq 0.50$ admission with $\ustored \geq 0.30$. \\
\bottomrule
\end{tabular}
\end{center}

\paragraph{Two-phase sweep protocol.}
Running all 14 training-time variants at the confirmatory
configuration ($n=5000$ SIL per cycle, 10 cycles, 3 seeds) exceeds the
project's self-funded compute envelope. The study therefore uses a
two-phase protocol.\footnote{The screening pass is a methodological
instrument to allocate confirmatory budget to the highest-$\Delta$CES
candidates; screening runs are not reported in the main ablation table.
Variant selection is pre-registered in
\texttt{caem-implementation-log.md} (Session~44) before the confirmatory
pass begins.} Phase~1 (self-funded, approximately USD~200) executes a
screening pass at $n=1500$ SIL for 3~cycles on a single seed (42) across
the entire 16-variant registry, ranks variants by point-estimate
$\Delta$CES, and then commits budget to the confirmatory pass at
$n=5000$ SIL for 10~cycles on a single seed for the top candidates
plus the reference run. Phase~2 (post-funding, approximately
+USD~700) re-runs the confirmatory pass for seeds 123 and 456 across
every cyclic variant, enabling mean$\pm$std reporting and
within-variant McNemar tests. Phase~1 results will be disclosed as
point estimates with an explicit single-seed limitation paragraph in
\S\ref{sec:ablation-results}; Phase~2 tables, once collected, will
supersede them with variance-bearing reports.

\paragraph{Metric reported.}
Each variant row reports the five CES axes
(ACC, EPI, RET, CAL, VER) and the geometric-mean scalar CES at the
final cycle of its own trajectory, plus $\Delta\text{CES}$ relative to
the reference run with independent-seeds standard-deviation propagation
once Phase~2 data is available. The per-cycle CES trace is also plotted
for the reference run and for the ablation variants whose $\Delta$CES
pattern across cycles is diagnostic (typically
\texttt{no\_self\_improvement}, \texttt{no\_verification}, and
\texttt{no\_retroverify}).

% ---------------------------------------------------------------
\section{Main Results}
\label{sec:main-results}
% ---------------------------------------------------------------

This section reports the seven metric-category tables introduced in
\cref{sec:metrics}, together with the three empirical validations of
the theoretical claims stated in~\cref{sec:theoretical-analysis}.
Each table below carries a caption and label already referenced from
the methodology paragraphs; numbers are populated from the
confirmatory-pass outputs at Cycle~$N$ and are reported with 95\%
confidence intervals where applicable.

% ---- Placeholder scaffolds for the seven metric-category tables ----
\begin{table}[t]
\centering
\caption{Headline accuracy and latency. Populated from
\texttt{outputs/eval/\{bm\}\_cycle\{N\}.json} via
\texttt{eval/reporting.build\_table\_headline}.}
\label{tab:headline}
\begin{tabular}{lrrrr}
\toprule
Benchmark & EM & $\mathrm{EM}_{T_1}$ & $\mathrm{EM}_{T_2}$ & $\mathrm{EM}_{T_3}$ \\
\midrule
\multicolumn{5}{c}{\emph{[To be populated from confirmatory-pass outputs.]}}\\
\bottomrule
\end{tabular}
\end{table}

\begin{table}[t]
\centering
\caption{Calibration: ECE, Brier score, and AUROC on pooled
$(\ustored, \mathrm{correct})$ pairs per cycle. Populated from
\texttt{eval/metrics.\{ece,brier,auroc\}}.}
\label{tab:calibration}
\begin{tabular}{lrrr}
\toprule
Cycle & ECE & Brier & AUROC \\
\midrule
\multicolumn{4}{c}{\emph{[To be populated from confirmatory-pass outputs.]}}\\
\bottomrule
\end{tabular}
\end{table}

\begin{table}[t]
\centering
\caption{Hallucination decomposition by sub-type
(confabulation, ungrounded, contradiction).}
\label{tab:halluc}
\begin{tabular}{lrrr}
\toprule
Cycle & Confabulation & Ungrounded & Contradiction \\
\midrule
\multicolumn{4}{c}{\emph{[To be populated from confirmatory-pass outputs.]}}\\
\bottomrule
\end{tabular}
\end{table}

\begin{table}[t]
\centering
\caption{Routing-distribution audit. Fraction of queries dispatched
to Tier~1, Tier~2, and Tier~3 per cycle, with the safety-override
fraction reported separately. Populated from
\texttt{outputs/eval/routing\_\{cycle\}.json}.}
\label{tab:routing}
\begin{tabular}{lrrrr}
\toprule
Cycle & Tier~1 & Tier~2 & Tier~3 & Override \\
\midrule
\multicolumn{5}{c}{\emph{[To be populated from confirmatory-pass outputs.]}}\\
\bottomrule
\end{tabular}
\end{table}

\begin{table}[t]
\centering
\caption{Verifier-signal attribution. Per-signal AUROC against
per-query correctness, reported for each of the nine signals in
the composite.}
\label{tab:verifier-attrib}
\begin{tabular}{lr}
\toprule
Signal & AUROC \\
\midrule
\multicolumn{2}{c}{\emph{[To be populated from confirmatory-pass outputs.]}}\\
\bottomrule
\end{tabular}
\end{table}

\begin{table}[t]
\centering
\caption{Episodic-memory growth and storage-outcome distribution by
cycle (\textsc{Store} / \textsc{Deferred} / \textsc{Abstain} /
\textsc{Discard}).}
\label{tab:memory-growth}
\begin{tabular}{lrrrrr}
\toprule
Cycle & $|\mathcal{M}|$ & Store & Deferred & Abstain & Discard \\
\midrule
\multicolumn{6}{c}{\emph{[To be populated from confirmatory-pass outputs.]}}\\
\bottomrule
\end{tabular}
\end{table}

\begin{table}[t]
\centering
\caption{CES composite and its five axes (ACC, EPI, RET, CAL, VER)
per cycle, with 95\% paired-seed confidence intervals.}
\label{tab:ces}
\begin{tabular}{lrrrrrr}
\toprule
Cycle & CES & ACC & EPI & RET & CAL & VER \\
\midrule
\multicolumn{7}{c}{\emph{[To be populated from confirmatory-pass outputs.]}}\\
\bottomrule
\end{tabular}
\end{table}

\begin{table}[t]
\centering
\caption{Grounding diagnostics per cycle. Unsupported-correct rate and
ungrounded-assertion rate, computed as defined in \cref{sec:metrics}.}
\label{tab:grounding}
\begin{tabular}{lrr}
\toprule
Cycle & Unsupported-correct (\%) & Ungrounded-assertion (\%) \\
\midrule
\multicolumn{3}{c}{\emph{[To be populated from confirmatory-pass outputs.]}}\\
\bottomrule
\end{tabular}
\end{table}

\begin{table}[t]
\centering
\caption{Continual-learning diagnostics: MMLU retention ratio,
backward transfer (BWT), and forward transfer (FWT) per cycle.}
\label{tab:continual}
\begin{tabular}{lrrr}
\toprule
Cycle & MMLU retention $\rho$ & BWT & FWT \\
\midrule
\multicolumn{4}{c}{\emph{[To be populated from confirmatory-pass outputs.]}}\\
\bottomrule
\end{tabular}
\end{table}

\begin{table}[t]
\centering
\caption{Data-purity validation (\cref{thm:purity}). Per-cycle
base accuracy $p$, verifier true-positive rate $\alpha$, theorem
prediction $P_{\text{theory}} = p\alpha / (p\alpha + (1-p)(1-\alpha))$,
and measured post-storage purity $P_{\text{obs}}$.}
\label{tab:purity}
\begin{tabular}{lrrrr}
\toprule
Cycle & $p$ & $\alpha$ & $P_{\text{theory}}$ & $P_{\text{obs}}$ \\
\midrule
\multicolumn{5}{c}{\emph{[To be populated from confirmatory-pass outputs.]}}\\
\bottomrule
\end{tabular}
\end{table}

\begin{table}[t]
\centering
\caption{Purity convergence trajectory (\cref{thm:convergence}).
Per-cycle observed purity $P_c$, first-difference increment
$\Delta P_c = P_c - P_{c-1}$, and the fitted envelope predicted by
the theorem's upper bound $p_{c+1} \le P_c$.}
\label{tab:purity-convergence}
\begin{tabular}{lrrr}
\toprule
Cycle & $P_c$ & $\Delta P_c$ & Envelope bound \\
\midrule
\multicolumn{4}{c}{\emph{[To be populated from confirmatory-pass outputs.]}}\\
\bottomrule
\end{tabular}
\end{table}

\begin{table}[t]
\centering
\caption{Statistical significance of \caem{} versus each external
baseline (B1--B7). Paired McNemar test with Edwards correction per
benchmark, bootstrap-resampled 95\% CI for EM difference.}
\label{tab:sig-test}
\begin{tabular}{llrrr}
\toprule
Baseline & Benchmark & $\Delta$EM & McNemar $p$ & 95\% CI \\
\midrule
\multicolumn{5}{c}{\emph{[To be populated from confirmatory-pass outputs.]}}\\
\bottomrule
\end{tabular}
\end{table}

\begin{figure}[t]
\centering
\fbox{\parbox{0.8\textwidth}{\centering\emph{[Reliability diagram
placeholder. Ten equal-width bins of $\ustored$ on the $x$-axis;
observed accuracy on the $y$-axis. Populated from the pooled
$(\ustored,\mathrm{correct})$ stream at Cycle~$N$.]}\\[4pt]\vspace{4em}}}
\caption{Reliability diagram on pooled $(\ustored, \mathrm{correct})$
pairs at Cycle~$N$, with ten equal-width confidence bins. Departure
from the $y=x$ identity line quantifies miscalibration; the ECE
value in Table~\ref{tab:calibration} is its weighted $L_1$ norm.}
\label{fig:reliability}
\end{figure}

\begin{figure}[t]
\centering
\fbox{\parbox{0.8\textwidth}{\centering\emph{[Purity trajectory
placeholder. Cycle index on the $x$-axis; measured pool purity
$P_c$ on the $y$-axis, with the theorem's predicted envelope
overlaid. Populated from \texttt{outputs/purity\_validation/
theory\_validation.json}.]}\\[4pt]\vspace{4em}}}
\caption{Pool-purity trajectory across cycles $0,\ldots,N$.
Solid line: measured $P_c$ from the purity-validation split; dashed
line: envelope predicted by \cref{thm:monotonicity} and
\cref{thm:convergence}.}
\label{fig:purity-trajectory}
\end{figure}

\paragraph{Theoretical validations.} Three empirical checks close the
chapter. First, the data-purity curve from~\cref{thm:purity} is plotted
against the measured Cycle~$t$ pool-purity trace. Second, the
monotonicity prediction of~\cref{thm:monotonicity} is checked by
verifying that the measured trace is non-decreasing in the
$\alpha > \tfrac{1}{2}$ regime. Third, the convergence claim
of~\cref{thm:convergence} is checked by fitting the late-cycle
increments $\Delta P_t$ against the predicted envelope.

% ---------------------------------------------------------------
\section{Ablation Results}
\label{sec:ablation-results}
% ---------------------------------------------------------------

This section reports the sixteen-variant ablation table introduced
in~\cref{sec:ablation}. For every variant the table records the CES
composite, its five axes, the $\Delta$CES against the full-CAEM
reference row, and the paired-seed significance test. Screening-pass
variants whose $\Delta$CES crosses the significance band are rerun
confirmatory and reported with the full five-axis decomposition;
remaining variants are reported with CES and $\Delta$CES only. The
table scaffold, metric columns, and row ordering are fixed in the
ablation registry (Table~\ref{tab:ablation-registry}); Phase~2
populates the numbers.

\begin{table}[t]
\centering
\caption{Ablation results. CES and five-axis decomposition for each
variant in the registry, with paired-seed $\Delta$CES against the
full-CAEM reference row.}
\label{tab:ablation-main}
\begin{tabular}{lrrrrrrr}
\toprule
Variant & CES & ACC & EPI & RET & CAL & VER & $\Delta$CES \\
\midrule
\multicolumn{8}{c}{\emph{[To be populated from confirmatory-pass outputs.]}}\\
\bottomrule
\end{tabular}
\end{table}

% ---------------------------------------------------------------
\section{Chapter Summary}
\label{sec:ch5-summary}
% ---------------------------------------------------------------

Chapter~5 fixes the experimental protocol for the Phase~2 run: a
seven-benchmark factual-QA panel with an MMLU retention control, the
CES geometric-mean composite with its five validated axes, the eight
external baselines (plus one optional STaR ceiling), and the
sixteen-variant cyclic ablation registry. The seven metric-category
tables (\cref{tab:headline,tab:calibration,tab:halluc,tab:routing,tab:verifier-attrib,tab:memory-growth,tab:ces})
and the ablation table (\cref{tab:ablation-main}) are scaffolded
above and populated by Phase~2; the three theoretical validations
(\cref{thm:purity,thm:monotonicity,thm:convergence}) are checked
against the same confirmatory-pass trace. The methodological claims
of~\cref{ch:requirements,ch:methodology} are the objects under test.
